% Versión para Overleaf: Punto 2 (Calidad del producto y del proceso)
\documentclass[a4paper,11pt]{article}
\usepackage[spanish]{babel}
\usepackage[utf8]{inputenc}
\usepackage[T1]{fontenc}
\usepackage{geometry}
\usepackage{graphicx}
\usepackage{booktabs}
\usepackage{tabularx}
\usepackage{caption}
\usepackage{hyperref}
\geometry{margin=2.2cm}

	itle{E4 — Evaluación de la calidad del producto y del proceso: medHistory}
\author{Claudia Mateo \and Yousra Jebari}
\date{}

\begin{document}
\maketitle

\section*{2. Evaluación de la calidad del producto y del proceso}

\subsection*{2.1 Calidad del producto}

Actualmente, \textbf{medHistory} se encuentra en un estado funcional parcialmente estable. Sus módulos principales están operativos, pero presenta carencias en validación y automatización que impiden considerarlo listo para un entorno de producción real.

A nivel arquitectónico se ha optado por una separación estricta entre frontend y backend mediante una API REST en PHP conectada a PostgreSQL. Como alternativa inicial se valoró el uso de Node.js; sin embargo, se descartó por la mayor curva de aprendizaje y el riesgo de incumplimiento de plazos. La elección de PHP se basa en la viabilidad técnica y rapidez de desarrollo, permitiendo integración progresiva.

\subsubsection*{Métricas y estado global del proyecto}

Para evaluar el estado actual se comparó esfuerzo planificado vs. esfuerzo real y se analizó el estado funcional del sistema. El avance funcional aproximado es del \textasciitilde65\%, con módulos principales implementados pero no completamente estabilizados. Funcionalidades clave (autenticación, perfil de salud y gestión de consultas) son operativas; otros módulos tienen limitaciones o están parcialmente implementados.

Principales observaciones:
\begin{itemize}
	\item El frontend se ha mantenido alineado con la planificación.
	\item El backend ha reducido parcialmente su alcance respecto a lo planificado.
	\item La fase de pruebas ha sido infraejecutada y constituye la principal limitación actual.
\end{itemize}

Impactos derivados:
\begin{itemize}
	\item \textbf{Tiempo}: aumento del tiempo de validación manual tras cada cambio, ralentizando la entrega.
	\item \textbf{Calidad}: cobertura de tests automáticos estancada en 40--60\% (unitarios e integración básica). Ausencia de E2E reduce la fiabilidad.
	\item \textbf{Riesgo}: mayor probabilidad de regresiones silenciosas al introducir cambios.
\end{itemize}

Decisión técnica prioritaria: estabilizar el sistema priorizando validación y testing frente a nuevas funcionalidades.

\subsubsection*{Estado funcional por módulos}

Resumen técnico (pruebas manuales y validación directa):

\paragraph{Autenticación y seguridad (JWT)} Login y registro operativos en entorno de pruebas; tokens de acceso y refresh implementados. Pendiente mejorar UX en mensajes de error y aplicar rotación/expiración más estricta de refresh tokens para reducir riesgos de sesión.

\paragraph{Perfil de salud} CRUD completado y persistencia en PostgreSQL funcionando. Formularios con validaciones básicas; falta validación en tiempo real (inline) y mejoras en usabilidad.

\paragraph{Módulo de consultas} Permite creación, visualización y asignación; en pruebas manuales el flujo se completa. Falta una suite de pruebas automáticas de integración y E2E para garantizar estabilidad tras cambios.

\paragraph{Antecedentes familiares} Flujo de datos completo y persistente. Problemas de interfaz: selector de fechas poco intuitivo y validación de formatos demasiado rígida.

\paragraph{Recordatorios} Primera versión funcional; requiere optimización de UX y corrección del tratamiento de zonas horarias (desfase temporal actual).

\paragraph{Logs de auditoría} Implementados y útiles para trazabilidad. En condiciones normales funcionan correctamente, pero en pruebas de estrés presentan inestabilidad ante alto volumen de peticiones, lo que limita la escalabilidad.

\subsubsection*{Tabla resumen por módulo}
\begin{table}[ht]
\centering
\caption{Estado por módulos — problema identificado e impacto técnico}
\begin{tabularx}{\textwidth}{l l X}
	oprule
Funcionalidad & Estado & Problema identificado / Decisión adoptada \\
\midrule
Autenticación (JWT) & Estable & UX deficiente en errores y falta de rotación estricta de refresh tokens. Decisión: priorizar feedback visual y mapear códigos HTTP a mensajes claros. \\
Perfil de salud & Funcional & Ausencia de validaciones inline y selector de fechas poco usable. Decisión: rediseñar validaciones antes de añadir vistas nuevas. \\
Consultas médicas & Funciona parcialmente & Falta de E2E automatizadas. Decisión: paralizar nuevas funciones y priorizar tests de integración. \\
Antecedentes & Funcional & Selector de fechas poco intuitivo y validación rígida. Decisión: reemplazar componente calendario y flexibilizar el parser de fechas. \\
Recordatorios & Parcial & Desfase por zonas horarias. Decisión: almacenar en UTC en backend y delegar conversión en frontend. \\
Logs de auditoría & Operativo (con riesgo) & Rendimiento inestable bajo alta concurrencia. Decisión: planificar migración a colas asíncronas (deuda técnica). \\
\bottomrule
\end{tabularx}
\end{table}

\subsubsection*{Control de incidencias y validación con usuarios}

Control interno: 20 incidencias resueltas y 5 abiertas; tiempo medio de resolución \textasciitilde48 h. Prueba de usabilidad con 3 usuarios externos: feedback positivo en navegación, pero fricciones en errores de login y en selección de horarios.

Decisión técnica: documentar fricciones como Issues formales en el repositorio y priorizar su mitigación.

\subsubsection*{Test de usuario y validación (5 participantes)}

Con el objetivo de evaluar la usabilidad en condiciones reales y validar el comportamiento de las funcionalidades, se realizó una prueba con cinco personas externas al equipo. Se empleó la técnica "Thinking Aloud" (pedir a los usuarios que verbalicen sus acciones y dudas). No se les ofreció ayuda; el moderador anotó tiempos, errores y comentarios.

Perfiles seleccionados:
\begin{itemize}
	\item Usuario 1: Estudiante universitario (uso habitual de aplicaciones digitales).
	\item Usuario 2: Persona adulta sin perfil técnico (bajo nivel digital).
	\item Usuario 3: Usuario intermedio con experiencia en aplicaciones web.
	\item Usuario 4: Familiar de paciente (perfil cercano al caso real de uso).
	\item Usuario 5: Usuario con perfil técnico básico (conocimientos informáticos generales).
\end{itemize}

Tareas planteadas a cada participante (sin ayuda):
\begin{enumerate}
	\item Registro en el sistema.
	\item Inicio de sesión.
	\item Acceso al perfil de salud.
	\item Introducción de datos en formularios (antecedentes / perfil).
	\item Creación de una consulta médica.
\end{enumerate}

Resultados (resumen):
\begin{table}[ht]
\centering
\caption{Resultados del test (5 usuarios)}
\begin{tabularx}{0.9\textwidth}{l l c c c}
\toprule
Usuario & Perfil & Completó flujo & Tiempo (s) & Errores / Satisfacción (1--5) \\
\midrule
Usuario 1 & Estudiante & Sí & 95 & 0 / 5 \\
Usuario 2 & Adulto no técnico & Parcial & 240 & 2 / 2.5 \\
Usuario 3 & Intermedio & Sí & 130 & 1 / 4 \\
Usuario 4 & Familiar & Sí & 160 & 1 / 3.5 \\
Usuario 5 & Técnico básico & Sí & 120 & 1 / 4 \\
\bottomrule
\end{tabularx}
\end{table}

Aspectos positivos detectados:
\begin{itemize}
	\item 4 de 5 usuarios completaron el flujo principal sin asistencia.
	\item Navegación general clara y estructura de menús comprensible.
	\item Funcionalidades principales (login, perfil, consultas) entendidas sin explicación previa en la mayoría de casos.
\end{itemize}

Problemas detectados (frecuencia entre paréntesis):
\begin{itemize}
	\item Mensajes de error de login poco claros, confusión al introducir credenciales (3/5).
	\item Dificultades con selector de fechas en formularios (2/5).
	\item Errores de formato en campos (fechas, numéricos) por falta de validación inline (3/5).
	\item Un usuario no percibió confirmación de guardado y abandonó parcialmente una tarea (1/5).
\end{itemize}

Errores detectados y soluciones propuestas:
\begin{table}[ht]
\centering
\caption{Errores detectados y soluciones}
\begin{tabularx}{\textwidth}{l l l X}
\toprule
Tarea & Problema & Usuario(s) & Solución propuesta \\
\midrule
Login & Mensajes de error genéricos & U1, U2, U3 & Mapear códigos HTTP a mensajes claros; destacar errores in-line. \\
Fecha en formularios & Selector poco intuitivo & U2, U4 & Reemplazar componente calendario por uno accesible y móvil-friendly. \\
Validación de campos & Formatos rígidos y sin feedback & U1, U3, U5 & Añadir validación en tiempo real y ejemplos de formato. \\
Confirmación de guardado & Falta de feedback visual & U2 & Insertar mensajes/alertas de éxito y estados de guardado. \\
\bottomrule
\end{tabularx}
\end{table}

Interpretación y decisiones técnicas adoptadas:
\begin{itemize}
	\item Priorizar la mejora del feedback en autenticación (mensajes claros y mapeo de códigos HTTP).
	\item Implementar validación inline en formularios y ejemplos de formato en campos sensibles.
	\item Sustituir/mejorar el selector de fechas por uno más usable y accesible.
	\item Añadir confirmaciones visuales al guardar para evitar incertidumbres.
\end{itemize}

Conclusión de la validación: la prueba confirma que el sistema es funcional y comprensible, pero presenta fallos de usabilidad y validación que deben resolverse antes de ampliar funcionalidades.

\subsubsection*{Evidencias en el repositorio}

Aspectos verificables en el repositorio:
\begin{itemize}
	\item Separación arquitectónica: directorios \texttt{backend/controllers} y \texttt{backend/dao}.
	\item Configuración frontend: \texttt{frontend/package.json} (Tailwind CSS).
	\item Corrección activa: \texttt{frontend/static/js/dependientes.js} (fix por asincronía).
	\item Pipeline básico en GitHub Actions: archivo \texttt{.github/workflows/ci.yml} (requiere ajustes para ejecutar correctamente el build). 
\end{itemize}

\subsubsection*{Diagnóstico global del producto}

El sistema cumple con el propósito principal y es funcional, pero arrastra deuda técnica en aseguramiento de calidad por priorizar desarrollo sobre validación automatizada. El riesgo principal es la estabilidad, escalabilidad y mantenibilidad a medio plazo.

\subsection*{2.2 Calidad del proceso}

A nivel metodológico, el proceso permitió construir la base del producto, pero muestra fallos organizativos que ralentizan la integración y fomentan deuda técnica.

\subsubsection*{Problemas metodológicos y su impacto}

\paragraph{Trabajo directo sobre \texttt{main}} No se utilizaron ramas feature para aislar desarrollos. Impacto: cambios incompletos o inestables llegaron a la rama principal, obligando a ciclos completos de validación manual y rompiendo temporalmente el trabajo de la otra integrante.

\paragraph{Falta de red de seguridad automatizada} El flujo de CI no bloquea la subida de código defectuoso. Impacto: desconfianza en el código integrado y sobreesfuerzo de testeo manual rutinario.

\subsubsection*{Plan de mejora del proceso (replanificación)}

Medidas innegociables para mitigar pérdida de velocidad y riesgo:
\begin{enumerate}
	\item \textbf{Modificación del flujo Git (Prioridad Alta)}: prohibición de commits directos a \texttt{main}; uso obligatorio de ramas \texttt{feature/xxx} y Pull Requests revisados por la otra integrante.
	\item \textbf{Bloqueo por CI (Prioridad Media)}: configurar GitHub Actions para que los PRs ejecuten los tests unitarios y bloqueen el merge si fallan.
	\item \textbf{Reasignación de esfuerzo}: pausar desarrollo de módulos secundarios (ej. chat) y dedicar horas a resolución de deuda técnica y estabilización.
\end{enumerate}

\subsection*{2.3 Estado detallado de las funcionalidades y evidencias}

Para ofrecer transparencia, se presenta un mapeo de funcionalidades planificadas, su estado real y la evidencia en el repositorio (controladores, DAOs, scripts y assets del frontend). Este desglose permite identificar lo terminado, lo parcial y lo no implementado.

\begin{itemize}
	\item \textbf{Autenticación (JWT)}: Funciona. Evidencia: \texttt{backend/src/controllers/AuthController.php}, \texttt{frontend/static/js/auth.js}, \texttt{static/login.html}, \texttt{static/register.html}. Pendiente: mensajes de error y rotación de refresh tokens.
	\item \textbf{Perfil de salud}: Funciona. Evidencia: \texttt{PerfilSaludController.php}, \texttt{PerfilSaludDAO.php} y formularios en frontend.
	\item \textbf{Antecedentes familiares}: Funciona. Evidencia: \texttt{AntecedentesController.php}, \texttt{AntecedentesFamiliaresDAO.php}, \texttt{frontend/static/js/antecedentes.js}. Pendiente: UX selector de fechas.
	\item \textbf{Módulo de consultas}: Parcial. Evidencia: \texttt{ConsultasController.php}, \texttt{ConsultaDAO.php}. Pendiente: pruebas E2E.
	\item \textbf{Gestión de pacientes dependientes}: Parcial. Evidencia: \texttt{DependienteController.php}, \texttt{DependienteDAO.php}, \texttt{frontend/static/js/dependientes.js}.
	\item \textbf{Recordatorios}: Parcial. Evidencia: \texttt{RecordatoriosController.php}, \texttt{RecordatorioDAO.php}, \texttt{database/corregir_recordatorios.sql}. Pendiente: frontend e integración de zonas horarias.
	\item \textbf{Módulo Administrador}: Parcial. Evidencia: \texttt{AdminController.php}, \texttt{static/admin.html}, \texttt{static/js/admin.js}.
	\item \textbf{Roles y permisos}: Parcial/Resuelto parcialmente. Evidencia: \texttt{MedicoDAO.php} y lógica asociada; INC-03 documentada.
	\item \textbf{Logs de auditoría}: Parcial/Riesgo. Evidencia: \texttt{AuditoriaDAO.php}; riesgo por alta concurrencia.
	\item \textbf{CI / Integración continua}: No completamente funcional. Evidencia: \texttt{.github/workflows/ci.yml} (referencia a comandos de build que requieren ajuste).
	\item \textbf{Pruebas automatizadas (Unit/E2E)}: No implementadas.
	\item \textbf{Módulo Citas / Chat}: No implementados.
\end{itemize}

\bigskip
\noindent\textbf{Nota:} Inserte 2--3 capturas de pantalla reales que muestren interfaces clave o mensajes de error para incluir evidencia visual en la versión final.

\end{document}
